\begin{pastexam}[label={px:ccu-112-csb-01}]{中正資工所乙組 112 學年度 第 1 題}
（單選並說明理由）

(6 points) How many function calls does \texttt{a(15)} produce if we use this recursive function
written in C? Please choose the correct answer and explain why.
\begin{minted}[fontsize=\small]{c}
int a(int n){
    if(n == 0)
        return 1;
    else
        return n * a(n-1);
}
\end{minted}
A.~14 \qquad B.~15 \qquad C.~16 \qquad D.~17
\end{pastexam}

\begin{pastsoln}
答案是 C，共 $16$ 次呼叫（把一開始的 \texttt{a(15)} 也算進去）。

\texttt{a(n)} 在 $n \ne 0$ 時呼叫一次 \texttt{a(n-1)}，$n = 0$ 時直接回傳，不再呼叫。 令 $C(n)$ 是
執行 \texttt{a(n)} 總共發生的呼叫次數（含它自己），照程式寫出遞迴式（資結書第~2.1~節）：
\[
C(0) = 1, \qquad C(n) = 1 + C(n - 1) \quad (n \ge 1).
\]
展開得 $C(n) = n + C(0) = n + 1$，所以 $C(15) = 16$。 直接列出來也一樣：
\[
\texttt{a(15)} \to \texttt{a(14)} \to \cdots \to \texttt{a(1)} \to \texttt{a(0)},
\]
參數從 $15$ 到 $0$，共 $16$ 個。 這 $16$ 層同時疊在呼叫堆疊上，\texttt{a(0)} 回傳 $1$ 之後才一層
一層往回乘（資結書第~6.4~節），最後得到 $15!$。

逐項檢查：
\begin{itemize}
\item A.~14：少算了 \texttt{a(15)} 本身和 \texttt{a(0)}，錯。
\item B.~15：只算了「遞迴呼叫」\texttt{a(14)} 到 \texttt{a(0)}，沒算最外層的 \texttt{a(15)}。 如果題目
的意思是「\texttt{a(15)} 又引發了幾次呼叫」，答案會是 15；但一般說「\texttt{a(15)} 產生幾次函數呼叫」
包含這一次，作答時寫出自己的算法即可。 依一般讀法，錯。
\item C.~16：$n = 15, 14, \ldots, 0$ 各一次，對。
\item D.~17：多算一次，\texttt{a(0)} 不會再呼叫 \texttt{a(-1)}，錯。
\end{itemize}
\end{pastsoln}
